\chapter*{}
\thispagestyle{empty}

\cleardoublepage
\thispagestyle{empty}

\begin{center}
{\large\bfseries Análisis de seguridad y optimización de generadores pseudoaleatorios basados en LFSR}\\
\end{center}
\begin{center}
Joaquín Sergio García Ibáñez\\
\end{center}

\noindent{\textbf{Palabras clave}: LFSR, Berlekamp-Massey, complejidad lineal, generadores de flujo, criptografía simétrica, tests NIST}\\

\vspace{0.7cm}
\noindent{\textbf{Resumen}}\\

Los registros de desplazamiento con retroalimentación lineal (LFSR) son la base de muchos
generadores de secuencias cifrantes usados en criptografía de flujo. Por sí solos presentan
vulnerabilidades graves, en particular frente al algoritmo de Berlekamp-Massey, que permite
reconstruir el registro completo a partir de un número pequeño de bits observados. Este trabajo
analiza cuatro generadores basados en LFSR ---Geffe, Shrinking, Self-Shrinking y Alternating
Step--- comparando su complejidad lineal y su comportamiento frente a los tests estadísticos
NIST SP 800-22. El objetivo es identificar cuáles ofrecen mejores garantías de seguridad y
entender por qué.

\cleardoublepage

\thispagestyle{empty}

\begin{center}
{\large\bfseries Security analysis and optimization of LFSR-based pseudorandom generators}\\
\end{center}
\begin{center}
Joaquín Sergio García Ibáñez\\
\end{center}

\noindent{\textbf{Keywords}: LFSR, Berlekamp-Massey, linear complexity, stream ciphers, symmetric cryptography, NIST tests}\\

\vspace{0.7cm}
\noindent{\textbf{Abstract}}\\

Linear Feedback Shift Registers (LFSRs) are the building block of many keystream generators
used in stream ciphers. On their own they have serious weaknesses, most notably against the
Berlekamp-Massey algorithm, which can reconstruct the register from a short observed sequence.
This thesis analyses four LFSR-based generators ---Geffe, Shrinking, Self-Shrinking and
Alternating Step--- comparing their linear complexity and statistical behaviour using the NIST
SP 800-22 test suite. The goal is to identify which ones provide stronger security guarantees
and to understand the reasons behind the differences.

\chapter*{}
\thispagestyle{empty}

\noindent\rule[-1ex]{\textwidth}{2pt}\\[4.5ex]

Yo, \textbf{Joaquín Sergio García Ibáñez}, alumno de la titulación Grado en Ingeniería
Informática de la \textbf{Escuela Técnica Superior de Ingenierías Informática y de
Telecomunicación de la Universidad de Granada}, autorizo la ubicación de la siguiente copia
de mi Trabajo Fin de Grado en la biblioteca del centro para que pueda ser consultada por las
personas que lo deseen.

\vspace{6cm}

\noindent Fdo: Joaquín Sergio García Ibáñez

\vspace{2cm}

\begin{flushright}
Granada a \myTime
\end{flushright}

\chapter*{}
\thispagestyle{empty}

\noindent\rule[-1ex]{\textwidth}{2pt}\\[4.5ex]

D. \textbf{Jesús García Miranda}, Profesor del Departamento de Álgebra de la Universidad de Granada.

\vspace{0.5cm}

\textbf{Informa:}

\vspace{0.5cm}

Que el presente trabajo, titulado \textit{\textbf{Análisis de seguridad y optimización de
generadores pseudoaleatorios basados en LFSR}}, ha sido realizado bajo su supervisión por
\textbf{Joaquín Sergio García Ibáñez}, y autoriza la defensa de dicho trabajo ante el tribunal
que corresponda.

\vspace{0.5cm}

Y para que conste, expide y firma el presente informe en Granada a \myTime

\vspace{1cm}

\textbf{El director:}

\vspace{5cm}

\noindent \textbf{Jesús García Miranda}

\chapter*{Agradecimientos}
\thispagestyle{empty}

\vspace{1cm}

